\section{Machine Learning}
\label{sec:machine-learning}

\subsection{Machine Learning Paradigms and Tasks}
\label{sec:machine-learning-paradigms-and-tasks}

Machine learning refers to algorithms that improve their performance 
at some task through experience 
\parencite{mitchell1997machine_learning}. 
Such algorithms can be broadly categorized according to the type of 
supervision they receive during that experience, also called training. 
There are three main paradigms commonly distinguished in the literature: 
supervised learning, unsupervised learning, and reinforcement learning 
\parencite{bishop2006pattern_recognition_machine_learning}
\parencite[Chapter~5]{Goodfellow-et-al-2016}
\parencite{sarker2021machine_learning_algorithms_real_world_applications_research_directions}.
The following sections outline each paradigm and introduce some common 
tasks.

In supervised learning, the training data consists of input vectors paired 
with corresponding target values. The two most common supervised tasks are 
classification, where the goal is to assign an input to one of a finite number 
of discrete categories, and regression, where the goal is to predict a continuous 
numerical value 
\parencite{bishop2006pattern_recognition_machine_learning}
\parencite[Chapter~5]{Goodfellow-et-al-2016}.

In unsupervised learning, no target values are provided. The algorithm must 
instead identify structure in the data directly 
\parencite{bishop2006pattern_recognition_machine_learning}. 
Common tasks include clustering, dimensionality reduction, and anomaly detection 
\parencite{sarker2021machine_learning_algorithms_real_world_applications_research_directions}
\parencite[Chapter~5]{Goodfellow-et-al-2016}.
Clustering groups data points such that objects within the same group are more 
similar to each other than to those in other groups 
\parencite{sarker2021machine_learning_algorithms_real_world_applications_research_directions}. 
Dimensionality reduction simplifies data by reducing the number of features, 
which leads to better human interpretations and lower computational costs
\parencite{sarker2021machine_learning_algorithms_real_world_applications_research_directions}. 
Anomaly detection identifies observations that deviate significantly from the 
expected distribution, and is used in applications such as fraud detection and 
fault monitoring 
\parencite[Chapter~5]{Goodfellow-et-al-2016}.
Semi-supervised learning occupies a middle ground, operating on both labeled 
and unlabeled data, which is useful in contexts where labeled examples are scarce 
\parencite{sarker2021machine_learning_algorithms_real_world_applications_research_directions}. 
It is worth noting that the boundary between supervised and unsupervised learning 
is not always well-defined, and many methods can be applied across both types 
of problems \parencite[Chapter~5]{Goodfellow-et-al-2016}.

Reinforcement learning differs fundamentally from the other paradigms. Rather 
than learning from a fixed dataset, an agent interacts with an environment and 
learns through trial and error, receiving rewards or penalties based on its actions 
\parencite{bishop2006pattern_recognition_machine_learning}
\parencite{sarker2021machine_learning_algorithms_real_world_applications_research_directions}. 
Within reinforcement learning, methods can be broadly divided according to what 
they learn. In value-based methods, a value function is learned from which the 
policy is derived either explicitly or implicitly, whereas in policy-based methods 
the policy is learned directly without the need for a value function 
\parencite{sewak2019policy_based_reinforcement_learning_approaches}
\parencite[Chapter~11]{sutton2018reinforcement_learning}. 
A third class, actor-critic methods, combines both approaches by maintaining 
a separate policy structure and value function simultaneously 
\parencite[Chapter~11]{sutton2018reinforcement_learning}.

\begin{figure}[H]
    \centering
    \includegraphics[width=\textwidth]{Figures/02-background/ML-Paradigms.png}
    \caption[ML paradigms]{Overview of machine learning paradigms with common tasks.}
    \label{fig:ml_paradigms}
\end{figure}

\subsection{Machine Learning Methods}
\label{sec:machine-learning-methods}

The following sections introduce a selection of common and most established machine learning 
methods and place them within the paradigms and tasks described above. 
The aim is not to provide a comprehensive description of each method or 
to cover underlying mathematical details, but to give an overview of 
common method types and what fundamentally distinguishes them. The section then gives a brief 
overview of attributes that are relevant to consider when selecting 
a method for a given problem.

\subsubsection*{Supervised methods}

Tree-based methods learn by recursively partitioning the input space 
according to feature values. A decision tree splits data into regions 
by asking a sequence of questions about the input features, assigning 
each region a predicted output. They can be used for both 
classification and regression tasks and require little data 
preparation, such as feature scaling 
\parencite[Chapter~6]{geron2019hands_on_machine_learning,
sarker2021machine_learning_algorithms_real_world_applications_research_directions}. 
Their decisions are intuitive and easy to interpret, as the 
classification rules they produce can even be applied manually 
\parencite[Chapter~6]{geron2019hands_on_machine_learning}.
Random forests extend this by training many trees on different 
subsets of the data and aggregating their predictions, which makes 
them more robust and generally improves predictive accuracy 
\parencite{breiman2001random_forests}. 
However, combining many trees makes the model harder to interpret 
than a single decision tree 
\parencite[Chapter~6]{geron2019hands_on_machine_learning}. 
Gradient boosting methods such as XGBoost build trees sequentially 
to improve on the predictions of previous trees, and are known for 
strong generalisation performance on large datasets
\parencite{chen2016xgboost,
sarker2021machine_learning_algorithms_real_world_applications_research_directions}.

Linear and polynomial regression predict a continuous output and 
are among the most widely used methods for regression tasks 
\parencite{sarker2021machine_learning_algorithms_real_world_applications_research_directions}. 
Regularised variants constrain the model weights to improve 
generalisation. Ridge regression keeps weights small but never 
eliminates them entirely, while LASSO regression tends to set some 
weights exactly to zero, which effectively performs feature 
selection and produces a sparse model 
\parencite{tibshirani1996regression_lasso_ridge, 
geron2019hands_on_machine_learning}. 
Logistic regression estimates the probability of class membership 
using a linear combination of input features, and is commonly used 
for classification tasks, particularly when a simple and 
interpretable model is preferred 
\parencite[Chapter~4]{geron2019hands_on_machine_learning,
sarker2021machine_learning_algorithms_real_world_applications_research_directions}.

Support vector machines find a decision boundary that maximises 
the margin between classes 
\parencite{cortes1995support_vector_networks}. 
By applying kernel functions, SVMs can model non-linear decision 
boundaries and are effective in high-dimensional spaces. They can 
be used for both classification and regression tasks 
\parencite[Chapter~5]{geron2019hands_on_machine_learning,
sarker2021machine_learning_algorithms_real_world_applications_research_directions}.

K-nearest neighbours classifies a new data point based on the 
majority class among its closest neighbours in the training data, 
using a distance measure such as Euclidean distance. The method 
requires no explicit training step but can be computationally 
expensive at prediction time for large datasets 
\parencite{cunningham2021k_nearest_neighbour_classifiers_tutorial,
sarker2021machine_learning_algorithms_real_world_applications_research_directions}.

\subsubsection*{Unsupervised methods}

Clustering algorithms group data points based on similarity without 
the use of labels. K-means partitions data into a fixed number of 
clusters by assigning each instance to the nearest centroid and 
updating the centroids iteratively until convergence 
\parencite[Chapter~9]{geron2019hands_on_machine_learning}. 
The number of clusters must be specified in advance, and the 
algorithm is sensitive to the initial centroid placement 
\parencite[Chapter~9]{geron2019hands_on_machine_learning}. 
DBSCAN takes a different approach by defining clusters as dense 
regions in the data space, which allows it to find clusters of 
arbitrary shape and identify outliers as noise without requiring 
the number of clusters to be specified in advance 
\parencite{ester1996dbscan,
sarker2021machine_learning_algorithms_real_world_applications_research_directions}.

For dimensionality reduction, principal component analysis (PCA) 
identifies the axes of greatest variance in the data and projects 
it onto a lower-dimensional space defined by these axes, called 
principal components 
\parencite{hotelling1932pca,
geron2019hands_on_machine_learning}. 
This reduces computational cost and can aid visualisation while 
preserving as much of the original variance as possible 
\parencite[Chapter~8]{geron2019hands_on_machine_learning,
sarker2021machine_learning_algorithms_real_world_applications_research_directions}.

Anomaly detection methods identify observations that deviate 
significantly from the expected distribution. Isolation Forest 
isolates anomalies by randomly partitioning the data, exploiting 
the fact that anomalies are few and structurally different from 
normal observations 
\parencite{liu2008isolation_forest}. 
One-class SVM, a variant of the support vector machine, learns a 
boundary around normal data and flags observations outside this 
boundary as anomalies 
\parencite[Chapter~9]{geron2019hands_on_machine_learning}.

\subsubsection*{Reinforcement learning methods}

Within reinforcement learning, methods differ primarily in what 
they learn. Value-based methods such as Q-learning estimate the 
expected return for each action in a given state and derive a 
policy from these estimates 
\parencite{watkins1992q_learning}. 
Policy-based methods such as policy gradient algorithms optimise 
the policy directly, which makes them better suited for problems 
with continuous or high-dimensional action spaces 
\parencite{sutton1999policy_gradient_methods_reinforcement_learning}. 
Actor-critic methods combine both approaches by maintaining 
a separate policy and value function simultaneously 
\parencite{zhang2020taxonomy_reinforcement_learning_algorithms}.

Figure~\ref{fig:ml_paradigms_methods} (on the following page) provides a structured overview 
of common machine learning methods organised by paradigm and task. 
It includes the methods described in the sections above, as well as 
additional commonly used methods not covered in detail here. 
The overview is based on a combination of taxonomies and categorisations described in the literature, 
and specifically draws on the works of
\textcite{sarker2021machine_learning_algorithms_real_world_applications_research_directions}, 
\textcite{geron2019hands_on_machine_learning}, and 
\textcite{zhang2020taxonomy_reinforcement_learning_algorithms}.

\begin{landscape}
\begin{figure}[H]
    \centering
    \includegraphics[width=1.0\linewidth, keepaspectratio]{Figures/02-background/ml-paradigms-methods-updated.png}
    \caption[ML methods by paradigm and task]{Overview of common machine learning methods categorised by
paradigms and tasks, based on
\textcite{sarker2021machine_learning_algorithms_real_world_applications_research_directions},
\textcite{geron2019hands_on_machine_learning}, and
\textcite{zhang2020taxonomy_reinforcement_learning_algorithms}.}
    \label{fig:ml_paradigms_methods}
\end{figure}
\end{landscape}

\subsection{Neural Networks and Deep Learning}
\label{sec:neural-networks-and-deep-learning}

Neural networks and deep learning are often treated as a separate category 
in the machine learning literature.
The following sections give a brief introduction to some of the most 
common neural network architectures and their typical use cases.

Neural networks are machine learning models inspired by the structure of the 
human brain, consisting of processing units organised in input, hidden, and 
output layers 
\parencite{shrestha2019review_deep_learning_algorithms_architectures}. 
They can be applied to a wide range of problems, including 
classification, regression, clustering, dimensionality 
reduction, computer vision, and natural language processing 
\parencite{shrestha2019review_deep_learning_algorithms_architectures}.
Deep learning refers to neural networks with many layers. By stacking 
multiple layers, the network learns increasingly abstract representations 
of the input data, where lower layers detect simple patterns and higher 
layers combine these into more complex concepts 
\parencite{lecun2015deep_learning}. 
This makes deep learning particularly well suited for complex, 
high-dimensional data such as images, audio, and text.
\parencite{lecun2015deep_learning, 
alzubaidi2021review_deep_learning_concepts}. 
A practical limitation is that deep learning models generally require 
large amounts of training data to perform well 
\parencite{alzubaidi2021review_deep_learning_concepts}.

\subsubsection*{Feedforward networks and MLP}

The simplest deep learning architecture is the multilayer perceptron (MLP), 
also known as the feedforward neural network. In an MLP, information flows 
in one direction from input to output through one or more hidden layers, 
without any feedback connections 
\parencite{shrestha2019review_deep_learning_algorithms_architectures, 
sarker2021machine_learning_algorithms_real_world_applications_research_directions}. 
Each node in one layer connects to each node in the following layer, and 
the network is trained using backpropagation to adjust the weights 
\parencite{lecun2015deep_learning}.

\subsubsection*{Convolutional neural networks}

Convolutional neural networks (CNNs) were originally developed 
for image recognition, where the spatial structure of the input 
is important \parencite{lecun2015deep_learning, lecun1998cnn}. 
Rather than connecting every neuron to all neurons in the previous layer, 
CNNs use local connections and shared weights, which reduces the number of 
parameters and makes the network faster to train 
\parencite{shrestha2019review_deep_learning_algorithms_architectures}. 
A typical CNN consists of convolutional layers, pooling layers, and fully 
connected layers. The convolutional layers progressively extract higher-level 
features from the input, while the pooling layers reduce the spatial size of 
the representations 
\parencite{pouyanfar2018survey_deep_learning_algorithms, 
sarker2021machine_learning_algorithms_real_world_applications_research_directions}. 
CNNs have achieved strong results in image recognition and computer vision, 
and are also applied to other data types such as audio and text 
\parencite{lecun2015deep_learning, pouyanfar2018survey_deep_learning_algorithms}.

\subsubsection*{Recurrent neural networks and LSTM}

Recurrent neural networks (RNNs) differ from feedforward networks in that 
the processing units form a cycle, allowing the network to maintain a state 
based on previous inputs 
\parencite{shrestha2019review_deep_learning_algorithms_architectures}. 
This makes RNNs well suited for sequential data such as time series, speech, 
and text 
\parencite{lecun2015deep_learning, pouyanfar2018survey_deep_learning_algorithms}. 
A known limitation of standard RNNs is the vanishing gradient problem, 
where information about earlier inputs is lost during training as the network 
processes longer sequences 
\parencite{pouyanfar2018survey_deep_learning_algorithms, 
shrestha2019review_deep_learning_algorithms_architectures}.

Long short-term memory networks (LSTMs) address this by introducing memory 
cells and gating mechanisms that control what information is stored, read, 
and written 
\parencite{hochreiter1997lstm, 
shrestha2019review_deep_learning_algorithms_architectures}. 
This allows LSTMs to learn from sequences with long-term dependencies, 
which makes them commonly used for time series analysis, speech recognition, 
and natural language processing 
\parencite{lecun2015deep_learning, 
sarker2021machine_learning_algorithms_real_world_applications_research_directions}.

\subsubsection*{Autoencoders}

Autoencoders are neural networks trained to compress input data into a 
lower-dimensional representation and then reconstruct the original input 
from this representation 
\parencite{shrestha2019review_deep_learning_algorithms_architectures}. 
Because they learn to reconstruct the input rather than predict a separate 
target, autoencoders are considered unsupervised models 
\parencite{shrestha2019review_deep_learning_algorithms_architectures}. 
They are commonly used for dimensionality reduction and anomaly detection, 
where observations that are difficult to reconstruct may indicate anomalies 
\parencite{sarker2021machine_learning_algorithms_real_world_applications_research_directions, 
shrestha2019review_deep_learning_algorithms_architectures}.

\begin{figure}[H]
    \centering
    \includegraphics[width=\textwidth]{Figures/02-background/neural-networks-table.pdf}
    \caption[Neural network architectures]{Summary and comparison of common neural network architectures, based on
    \textcite{sarker2021machine_learning_algorithms_real_world_applications_research_directions},
    \textcite{shrestha2019review_deep_learning_algorithms_architectures},
    \textcite{pouyanfar2018survey_deep_learning_algorithms}, and
    \textcite{lecun2015deep_learning}.}
    \label{fig:neural_networks_table}
\end{figure}

\subsubsection*{Deep reinforcement learning}

Deep reinforcement learning (DRL) combines deep learning with reinforcement 
learning to build agents that can operate in high-dimensional environments. 
The main motivation for using deep networks in RL is to handle complex, 
high-dimensional state spaces that cannot be represented by simple tabular 
or linear methods 
\parencite{zhang2020taxonomy_reinforcement_learning_algorithms}. 
A key development in DRL was the deep Q-network (DQN), which used a neural
network to approximate the Q-value function and enabled agents to learn 
directly from raw inputs such as image pixels 
\parencite{zhang2020taxonomy_reinforcement_learning_algorithms}. 
DRL has demonstrated strong performance in domains such as game playing, 
robotics, and control, where the state and action spaces are too large 
for traditional RL methods 
\parencite{shrestha2019review_deep_learning_algorithms_architectures, 
zhang2020taxonomy_reinforcement_learning_algorithms}.

\subsection{Attributes Relevant for Method Selection (MISSING)}
\label{sec:attributes-relevant-for-method-selection}

Attributes relevant for method selection include:
\begin{itemize}
    \item Data type: tabular, time series, image, text
    \item Data availability: large vs. limited data requirements
    \item Interpretability: interpretable vs. black-box methods
    \item Performance preferences: accuracy, training time, computational cost, prediction speed
\end{itemize}


\subsection{Approaches to Method Selection}
\label{sec:approaches-to-method-selection}

Selecting an appropriate machine learning method is not straightforward. 
As discussed in Section~\ref{sec:motivation}, the large number of
available algorithms, each with distinct strengths and limitations, creates 
a real challenge for practitioners. This challenge is formally described by 
the Algorithm Selection Problem \parencite{smith_miles2009algorithm_selection}, 
which asks which algorithm is likely to perform best for a given problem. 
The No Free Lunch theorems provide a theoretical basis for this difficulty, 
showing that no single algorithm performs best across all problems and that 
good selection requires understanding the characteristics of the problem at 
hand \parencite{wolpert1997no_free_lunch}.

Several practical tools have been developed to help with method selection. 
Scikit-learn provides a publicly available decision flowchart that guides 
users through questions about data size and task type to suggest candidate 
algorithms \parencite{sklearn_algorithm_cheatsheet}, and Microsoft Azure 
offers a similar cheat sheet for its machine learning platform 
\parencite{azure_algorithm_cheatsheet}. While these tools are accessible 
and widely used, they have clear limitations. 
The selection criteria used in both tools are mainly accuracy and prediction 
speed, something that can in practice lead to poor results \parencite{sala2018selecting_ml_algorithm}. 
Factors such as interpretability, data characteristics, and domain context 
are not considered, nor do they offer any guidance on how the problem 
itself should be formulated \parencite{chen2023ml_manufacturing_industry4}.

More structured frameworks have been proposed in the academic literature. 
\textcite{sala2018selecting_ml_algorithm} present a two-layer selection 
framework that incorporates algorithm characteristics such as scalability 
and robustness to outliers, in addition to learning type and response type. 
Domain-specific frameworks have also been proposed, such as that of 
\textcite{arena2024conceptual_framework_ml_algorithm_selection_predictive_maintenance}, 
who develop a structured set of decision variables for algorithm selection 
in predictive maintenance. However, as the authors note, most existing 
approaches focus narrowly on accuracy and computational requirements, 
and none provides a tool that connects algorithmic choices to the broader 
context of the application domain.

A related class of approaches is automated machine learning (AutoML), 
which seeks to automatically select, compose, and parametrise machine 
learning models to achieve optimal performance on a given task or dataset 
\parencite{waring2020automated_machine_learning_review_state_of_the_art_opportunities_healthcare}. 
AutoML systems automate parts of the machine learning pipeline, including 
data preparation, feature engineering, hyperparameter optimisation, and 
model generation, with the goal of reducing the demand for expert knowledge 
and making machine learning more accessible to domain experts 
\parencite{he2021automl_survey_state_of_the_art}. 
Within AutoML, the algorithm selection problem is typically treated as an
optimisation task, where dataset characteristics are used to predict which 
algorithm is likely to perform best 
\parencite{barbudo2023eight_years_automl_categorisation_review_trends}.

However, AutoML approaches have notable limitations in the context of 
decision support for non-experts. Most systems treat the selection and 
configuration process as a black box, providing little explanation of why 
a particular configuration was chosen 
\parencite{barbudo2023eight_years_automl_categorisation_review_trends}. 
This lack of transparency is a significant barrier to adoption, as users 
who do not understand or trust the output of a system are unlikely to apply 
it in practice 
\parencite{waring2020automated_machine_learning_review_state_of_the_art_opportunities_healthcare}. 
Furthermore, AutoML primarily addresses the later stages of the machine 
learning pipeline, and the phases considered inherently human, such as 
problem formulation and interpretation of results, have received little 
attention \parencite{barbudo2023eight_years_automl_categorisation_review_trends}. 
This means that AutoML does not address the earlier challenge of helping 
practitioners understand which type of problem they have and which general 
approach is appropriate. This gap motivates the 
work in this thesis.

\textbf{Burde siste setning om "this gap" stå her? Eller ta med noe lignende i motivasjon? Eller ingen av delene?}
